\section{Related Work}
\label{sec:related}
\subsubsection*{Distributed stream processing systems}
Apache Storm~\cite{toshniwal2014storm} is among the first distributed stream processing systems. It adopts the execution model in 
which the query is represented as a DAG of operators, and data flows from one operator to the next. Subsequent systems 
follow the same model and add more features to the original Storm design. Spark Streaming \cite{zaharia2013discretized}, 
Flink \cite{carbone2015apache} and Samza \cite{noghabi2017samza} support both batch and stream processing in a single framework and include
many stateful operators such as joins and aggregates. These systems handle operator failures, which can affect query 
correctness, by replicating the operators, checkpointing \cite{murray2013naiad}, or a combination thereof \cite{lin2016streamscope}. Other works focus on improving the performance of DAG 
execution by detecting bottlenecks via backpressure \cite{kulkarni2015twitter}, or by dynamically replicating the
operators to adapt to workload changes \cite{castro2013integrating,wu2015chronostream}. Our work demonstrates attacks 
against a generic DSP system, which can be applied to systems other than Flink.  

\subsubsection*{Outsourced data processing}
Designing an outsourced database running on an untrusted cloud requires making trade-offs between security, 
functionality, and performance. Cryptographic solutions such as CryptDB \cite{popa2011cryptdb} and 
Arx \cite{poddar2016arx} protect data confidentiality using cryptographic techniques, namely, using homomorphic and
functional encryption schemes. However, they support limited computation over the data, and have been shown to leak data 
in practical settings \cite{naveed2015inference}. Secure databases built on top of trusted hardware include 
Cipherbase~\cite{arasu2013orthogonal} and TrustedDB~\cite{bajaj2011trusteddb} that use older hardware.                      
More recent systems based on Intel SGX, such as ObliDB \cite{oblidb}, Opaque \cite{zheng2017opaque}, Obaldi \cite{crooks2018obladi}, focus on obliviousness to prevent leakage through memory access patterns \cite{priebe2018enclavedb}. They achieve high performance and more functionalities than purely cryptographic solutions, but their security assumptions are stronger, i.e. the hardware is secure. Furthermore, they incur non-trivial
overhead due to obliviousness.

Our work focuses data stream processing systems --- one type of database that handles continuous, infinite streams of 
data. In these systems, query confidentiality, which is overlooked in the other works targeting relational databsaes, is 
as important as data confidentiality. StreamBox-TZ~\cite{234964} and SecureStream~\cite{DBLP:journals/corr/abs-1805-01752} are the two most 
relevant streaming systems using secure hardware, but they also do not provide query confidentiality. 

\subsubsection*{Side channels}
Trusted hardware, in particular Intel SGX, has been a popular target for side channel attacks. The attacker gains
knowledge of the code and data running inside the enclaves via both hardware and software side channels. Since SGX
assumes a strong threat model, the attacker can manipulate page faults~\cite{van2017telling}, timer                         interrupt~\cite{van2017sgx}, and exception~\cite{cui2021smashex} to leak content from inside the enclave. Despite the           closed design, recent works have uncovered multiple hardware, micro-architectural side channels, including hardware              
cache~\cite{brasser2017software, van2021cacheout}, transient execution~\cite{van2018foreshadow}, and branch
prediction~\cite{lee2017inferring}. A comprehensive list of attacks against SGX can be found in~\cite{van2024sok}. Our 
work demonstrates another side channel specific to distributed stream processing systems. It is simpler to execute than 
other state-of-the-art channels, and we demonstrate that it is highly effective at leaking the application queries. How 
our attack can be improved with the other side channels is left as future work.